\chapter{Adverse Effects From Medicines}
\index{Adverse Effects From Medicines}

\section*{Definition and Overview}
An adverse effect, also called a side effect or adverse drug reaction, is an unwanted or harmful effect caused by a medicine when it is used at its normal dose. Adverse effects range from mild and annoying, such as nausea or drowsiness, to severe and life-threatening, such as a serious allergic reaction or organ damage. They are distinct from accidental overdose, drug interactions, and withdrawal symptoms, although all of these can be mistaken for one another.

Any medicine can cause adverse effects, including those bought without a prescription, complementary and herbal products, and vaccinations. Understanding which effects are common and expected, which need urgent attention, and how to report unexpected reactions is an important part of using medicines safely.

\section*{Epidemiology}
Adverse effects from medicines are very common. Roughly half of adults take at least one prescribed medicine, and older people often take five or more, which substantially increases the risk of interactions and harm. Adverse reactions are thought to contribute to up to 1 in 10 hospital admissions among older adults, and many of these are preventable.

Certain medicine groups account for most serious reactions, including blood thinners, diabetes medicines, blood-pressure medicines, antibiotics, and strong pain relievers. Most adverse effects are predictable from the way the medicine works (dose-related); true allergic reactions are rarer and occur only in a small minority of people.

\section*{Aetiology and Pathophysiology}
Adverse effects occur through several distinct mechanisms.

\begin{itemize}
    \item \textbf{Type A (dose-related):} predictable effects that are an extension of the medicine's normal action -- for example, bleeding with blood thinners, low blood sugar with insulin, or drowsiness with some antihistamines
    \item \textbf{Type B (idiosyncratic or allergic):} unpredictable reactions not related to dose, including drug allergies and rare severe reactions such as serious skin reactions or liver injury
    \item \textbf{Drug interactions:} another medicine, food, or herbal product alters how the first medicine works -- for example, grapefruit juice, alcohol, or St John's Wort
    \item \textbf{Medication errors:} taking the wrong dose, the wrong medicine, or a medicine at the wrong time
    \item \textbf{Individual factors:} age, kidney and liver function, pregnancy, and genetic variation all influence how the body handles a medicine
\end{itemize}

\section*{Clinical Features}
The features depend on the medicine, but some patterns are common.

\begin{itemize}
    \item \textbf{Common mild effects:} nausea, vomiting, diarrhoea, headache, dizziness, drowsiness, dry mouth, and rashes
    \item \textbf{Allergic reactions:} hives, itching, swelling of the lips or face, and in severe cases difficulty breathing and collapse (anaphylaxis)
    \item \textbf{Serious organ effects:} yellowing of the skin or eyes (liver), reduced urine output (kidneys), and unusual bruising, bleeding, or unexplained infections (blood cells)
    \item \textbf{Severe skin reactions:} widespread, blistering, or painful skin involvement
    \item \textbf{Late effects:} some reactions develop only after weeks or months of treatment, and some emerge when a medicine is stopped or reduced (withdrawal)
\end{itemize}

\section*{Differential Diagnosis}
When a new symptom develops in someone taking a medicine, the cause is not always the medicine.

\begin{itemize}
    \item An adverse effect versus progression of the condition being treated
    \item An adverse effect versus a new, unrelated illness
    \item A drug interaction versus an independent event
    \item An allergic reaction versus a non-allergic reaction to the medicine
    \item Withdrawal symptoms versus the treated condition returning
    \item Effects of diet, stress, or other substances
\end{itemize}

\section*{When to Seek Medical Care}
Many mild side effects settle on their own or can be managed with advice from a pharmacist. Seek medical advice promptly for persistent, severe, or worrying effects, and do not stop some medicines abruptly -- such as corticosteroids, anticonvulsants, or antidepressants -- without medical guidance, because sudden withdrawal can be dangerous.

\textbf{Contact your local emergency services for:}
\begin{itemize}
    \item Difficulty breathing, swelling of the face, lips, or throat, or collapse -- signs of a severe allergic reaction
    \item Widespread, blistering, or painful skin reactions
    \item Suspected overdose or deliberate self-harm with medicines -- contact local emergency services or the poison information service for urgent advice
    \item Chest pain, unusual bleeding, severe abdominal pain, or fainting
\end{itemize}

\section*{Diagnosis and Investigations}
Identifying an adverse effect relies on a careful review of the medicine history and the timing of symptoms.

\begin{itemize}
    \item \textbf{Medication review:} a complete list of all prescribed, over-the-counter, herbal, and recreational substances, with doses and dates started
    \item \textbf{Timing analysis:} the relationship between starting or changing a medicine and the onset of symptoms
    \item \textbf{Dechallenge and rechallenge:} observing whether symptoms improve when the medicine is stopped and return when it is restarted, under medical supervision
    \item \textbf{Blood tests:} liver, kidney, and full blood count, plus drug levels where relevant (for example, warfarin or lithium)
    \item \textbf{Allergy testing:} skin or blood tests for a suspected drug allergy, arranged through a specialist
    \item \textbf{Reporting:} serious or unexpected reactions can be reported through the medicines-safety system in your jurisdiction
\end{itemize}

\section*{Management}
Most adverse effects are managed simply, but serious ones require urgent care.

\begin{itemize}
    \item \textbf{Tell someone:} report the effect to a doctor or pharmacist promptly
    \item \textbf{Do not just stop:} stopping abruptly is dangerous for some medicines, so seek advice first unless the reaction is severe
    \item \textbf{Adjust the medicine:} a doctor may reduce the dose, change the timing, or switch to an alternative
    \item \textbf{Manage the symptom:} for example, taking the medicine with food for nausea, or taking it at night if it causes drowsiness
    \item \textbf{Treat severe reactions urgently:} stop the medicine and seek emergency care for anaphylaxis or serious organ reactions
    \item \textbf{Keep a record:} document the reaction in the medical record so the medicine is not prescribed again
\end{itemize}

\section*{Complications}
\begin{itemize}
    \item Anaphylaxis -- a life-threatening emergency
    \item Organ damage, including injury to the liver, kidneys, heart, or bone marrow
    \item Internal bleeding, particularly from blood-thinning medicines
    \item Falls and injuries in older people from dizziness or low blood pressure
    \item Confusion or delirium, especially in elderly patients
    \item Treatment failure if a medicine is stopped because of an unrecognised interaction
    \item Hospital admission and, rarely, death from a serious reaction
\end{itemize}

\section*{Prognosis and Outcomes}
Most adverse effects are mild and settle completely when the dose is adjusted or the medicine is stopped. Serious reactions are uncommon but can be dangerous if not recognised early. The risk of harm is greatly reduced when medicines are reviewed regularly, allergies and past reactions are documented, and patients understand what they are taking. A documented drug allergy should be treated as lifelong unless proven otherwise by a specialist.

\section*{Prevention and Self-Care}
\begin{itemize}
    \item \textbf{Medicine list:} keep an up-to-date list of all medicines, doses, allergies, and past reactions, and take it to every appointment
    \item \textbf{One pharmacy:} use a single pharmacy where possible so interactions can be checked
    \item \textbf{Read the information:} review the Consumer Medicines Information for each medicine
    \item \textbf{Ask questions:} ask your doctor or pharmacist which side effects to expect and which need urgent attention
    \item \textbf{Never share:} do not take another person's prescription medicines
    \item \textbf{Measure carefully:} measure liquid medicines with the proper device, not a kitchen spoon
    \item \textbf{Medicines review:} have a regular review with your GP or pharmacist, especially if you take several medicines
    \item \textbf{Medical alert:} wear a medical alert bracelet if you have a serious drug allergy
\end{itemize}

\section*{Sources and Further Reading}
The following sources provide reliable, up-to-date information on adverse effects of medicines.

\begin{itemize}
    \item NHS. \textit{Medicines A to Z: side effects of medicines and how to report them.} https://www.nhs.uk
    \item Mayo Clinic. \textit{Drug side effects: what to do when they occur.} https://www.mayoclinic.org
    \item MedlinePlus. \textit{Drug side effects and drug interactions.} https://medlineplus.gov
    \item World Health Organization. \textit{Pharmacovigilance: safety of medicines and adverse reactions.} https://www.who.int
    \item MSD Manual. \textit{Overview of drug reactions and interactions.} https://www.msdmanuals.com
    \item Centers for Disease Control and Prevention. \textit{Medication safety and adverse drug event prevention.} https://www.cdc.gov
\end{itemize}
